\section{Conclusion}
\label{sec:conclusion}
\enlargethispage{2pt}

This paper introduced the Tucker polynomial, which extends the linear map from
a Tucker core to the data tensor with $p$-fold interactions. Sharing the mode
factors across polynomial orders yields NTD-PL. This compact form reuses one
core and one set of factors, and it is also a low-rank Tucker tensor followed by
an entrywise polynomial response. Other finite scalar response bases can be used
with the same estimator when different response bases are appropriate.

The analysis established the relation between the general $p$-fold terms and
their shared-factor form, quantified the resulting rank and parameter cost, and
studied response approximation, nonconvex descent, and computational scaling.
The estimator combines a small ridge update for the response coefficients with
gradient-based Tucker updates. The diagnostic measures how much of a fixed
Tucker residual can be removed before joint NTD-PL training.

Controlled experiments show that NTD-PL remains close to Tucker for linear
measurements and separates as nonlinear energy increases. Hyperspectral
reconstruction and completion show improvements in RMSE and spectral angle at
matched Tucker rank, while higher-rank Tucker comparisons quantify the extra
parameters needed to reach similar RMSE. The mechanism experiments connect the
gain to the effective-rank increase after the response, and the cross-domain
results show that larger diagnostic scores are followed by larger NTD-PL gains.

Several directions remain open. Monotone splines or physics-informed response
families may be preferable when calibration constraints are known. Sharper
sample-complexity or identifiability results, richer noise models, and
domain-specific response priors are also natural extensions of the framework.
